% SGA 3, Expose VI_A, Loop 1 English reconstruction.
% Scope: paragraph 3.2.1 in full; paragraph 3.2.2 is excluded.
% Authority: Polo--Gille Exp6A-13oct24.pdf, component-local pages 22--23
% (combined-reader pages 325--326; printed pages 314--315).
% OCR and the Reinhold English Markdown are comparison witnesses only.
% The diagram is an exact 180-dpi crop of the controlling PDF:
% local p. 23, x=490,y=440,w=580,h=170 on a 1488x2105 render.

\medskip
\phantomsection\label{par:sga3-VIA-3.2.1}%
\noindent\textbf{3.2.1.---}
暂且假设 \(G\) 的每个点有饱和开邻域 \(W\),
使 \(G_*\) 在 \(W\) 上诱导的群胚有拟截面 (参见第 V 篇, \S~6).
则由\SGARefLink{lem:sga3-V-6.1}{第 V 篇, 6.1},
断言\SGARefLink{lem:sga3-V-6.1:item:i}{(i)}, (ii), (iii) 及 (iv)(a) 成立,
且 \(F\backslash G\) 在 \(A\) 上有限型.\begingroup
\renewcommand{\thefootnote}{\fnsymbol{footnote}}%
\footnote[1]{英译者注: 法文底本在此处以及\hyperref[note:sga3-VIA-39]{编者注 39} 中均印作 \(k\).
本部分证明的底为 Artin 局部环 \(A\), 所需陈述是``在 \(A\) 上有限型''.}%
\endgroup
此外, 在 (iv) 的假设下, 由于
\[
  G\longrightarrow F\backslash G
\]
忠实平坦且局部有限呈示, 断言 (b) 由 EGA~IV, 2.2.14 和 17.7.7 得到.
由第 IV 篇, 3.4.3.1, 5.2.2 和 5.1.6,
断言 (iv)(a\(^{\prime}\)) 由 (iv)(a) 得到.
最后, 断言 (iv)(c) 将在\hyperref[par:sga3-VIA-3.2.5]{3.2.5} 中证明,
而 (ii\(^{\prime}\)) 和 (iv)(d) 将在\hyperref[sec:sga3-VIA-5]{第~5~节}证明.

现在证明下述断言:
\[
  \tag{\(\dagger\)}\label{eq:sga3-VIA-3.2.1-dagger}
  \text{\(F\backslash G\) 的每个有限子集都包含于一个仿射开子集中.}
\]%
\footnote{\label{note:sga3-VIA-35}编者注: 指出如下事实: 若 \(A=k\) 为域,
则 \(F\backslash G\) 的每个拟紧开子集都拟射影
(对平滑代数群, 此结果由 Chow 给出);
参见 [Ray70], VI, \SGARefLink{cor:sga3-VIB-2.6}{2.6}.
而在 Artin 局部环 \(A=\mathbb C[\varepsilon]/(\varepsilon^2)\) 上,
存在并不射影的阿贝尔 \(A\)-概形 \(G\) (同上, XII, 4.2).}

若 \(U\) 为 \(G_*\) 在 \(G\) 的饱和开子集 \(W\) 上诱导的群胚的拟截面,
则 \(U\otimes_AA'\) 是 \(G_*\otimes_AA'\) 在 \(W\otimes_AA'\) 上
诱导的 \((\mathrm{Sch}/A')\)-群胚的拟截面.
此外, 若 \(U_*\) 为 \(G_*\) 在 \(U\) 上诱导的 \((\mathrm{Sch}/A)\)-群胚,
则 \(U_*\otimes_AA'\) 可等同于 \(G_*\otimes_AA'\)
在 \(U\otimes_AA'\) 上诱导的 \((\mathrm{Sch}/A')\)-群胚.%
\footnote{\label{note:sga3-VIA-36}编者注: 上述断言对每个基变换 \(A\to A'\) 都成立.}

由第 V 篇的证明可知, 商
\[
  X=F\backslash G
\]
的构造与这里考虑的这类基扩张 \(A\to A'\) 可交换.%
\footnote{\label{note:sga3-VIA-37}编者注: 下文\SGARefLink{par:sga3-VIA-4.6}{4.6} 将详细解释这一点.
须证明沿态射 \(p\), \(\lambda\) 和 \(\operatorname{pr}_2\) 取正像的构造
与平坦基变换 \(A\to A'\) 可交换.
由于 \(F\) 和 \(G\) 在 Artin 环 \(A\) 上有限型,
所涉及的态射 \(f\) 都拟紧且拟分离, 故
\[
  f_*(\mathcal O_X)\otimes_AA'
  =
  f'_*(\mathcal O_{X'})
\]
(采用显然的记号) 由 EGA~IV\textsubscript{1}, 1.7.21 得到.}

设 \(x_1,\ldots,x_n\) 为 \(X=F\backslash G\) 的点, 可假设它们为闭点,%
\footnote{\label{note:sga3-VIA-38}编者注: 事实上, 设 \(y_1,\ldots,y_n\) 为 \(X\) 的任意点.
由于 \(X\) 在 \(A\) 上有限型, 每个 \(y_i\) 的闭包包含闭点 \(x_i\),
且包含 \(x_i\) 的每个开子集都包含 \(y_i\).}
并设 \(g_1,\ldots,g_n\) 为 \(G\) 的闭点, 映到 \(x_1,\ldots,x_n\).
设 \(V\) 为 \(X\) 的处处稠密仿射开子集,%
\footnote{\label{note:sga3-VIA-39}编者注: 这样的开子集存在, 因为 \(X\) 在 \(A\) 上有限型:
\(X\) 只有有限多个不可约分支 \(C_1,\ldots,C_p\),
只需对每个 \(i\) 选取一个包含于下列集合的非空仿射开子集:
\[
  C_i-\bigcup_{j\ne i}C_j.
\]
(此外, 这里由 (ii\(^{\prime}\)) 可知各 \(C_i\) 不相交.)}
并设 \(U\) 为其在 \(G\) 中的逆像.
由\SGARefLink{prop:sga3-VIA-1.2}{1.2}, 存在为局部环且在 \(A\) 上有限自由的
\(A\)-代数 \(A'\), 使
\[
  G'=G\otimes_AA'
\]
中位于 \(g_1,\ldots,g_n\) 上方的点 \(g'_1,\ldots,g'_p\)
在 \(A'\) 上严格有理.%
\footnote{\label{note:sga3-VIA-40}编者注: 下文对原版作了扩充.}
由于态射 \(G'\to G\) 和 \(G\to X\) 为开态射,
\[
  U'=U\otimes_AA'
\]
在 \(G'\) 中稠密. 故开子集
\[
  \bigcap_{i=1}^p (U')^{-1}\cdot g'_i
\]
非空, 从而包含闭点 \(x\).
由\hyperref[prop:sga3-VIA-1.2]{1.2} (及\hyperref[sga3:VIA:0.4.1]{0.4.1}),
必要时扩大 \(A'\), 可假设 \(x\) 在 \(A'\) 上严格有理.
由于 \(x\in(U')^{-1}\cdot g'_i\), 有
\[
  g'_i\in U'\cdot x
  \qquad(1\leq i\leq p).
\]

用 \(V'\) 表示 \(V\) 在下列概形中的逆像:
\[
  X'=X\otimes_AA'.
\]
它是 \(X'\) 的仿射开子集, 也是 \(U'\) 在投影 \(G'\to X'\) 下的像.
由于右平移 \(r_x\) 是群胚 \(G_*\otimes_AA'\) 的自同构,
它诱导商 \(X'\) 的一个仍记为 \(r_x\) 的自同构. 因此,
\[
  V'\cdot x=r_x(V'),
\]
即 \(U'\cdot x\) 在 \(X'\) 中的像, 是 \(X'\) 的仿射开子集,
包含 \(g'_1,\ldots,g'_p\) 的像 \(x'_1,\ldots,x'_p\).

考虑投影 \(X\otimes_AA'\to X\)
在 \(X'=X\otimes_AA'\) 上定义的等价关系:
\begin{center}
  \phantomsection\label{diag:sga3-VIA-3.2.1-descent-equivalence}%
  \begin{tikzcd}[column sep=large]
    X\otimes_A A'\otimes_A A'
      \arrow[r,shift left=1.2ex,"d_1"]
      \arrow[r,shift right=1.2ex,"d_0"']
    &
    X\otimes_A A'
      \arrow[r]
    &
    X
  \end{tikzcd},
\end{center}
这里 \(d_0\) 和 \(d_1\) 由 \(A'\) 到 \(A'\otimes_AA'\) 的两个典范单射诱导.
由于 \(A'\) 是有限自由 \(A\)-代数, 设秩为 \(n\),
则 \(d_0\) 和 \(d_1\) 为秩 \(n\) 的有限局部自由态射.
因此可应用\SGARefLink{par:sga3-V-5-b}{第 V 篇, 5.b} 的论证,
它取自 SGA~1, VIII.7.6 的证明.
从而 \(x'_1,\ldots,x'_p\) 包含于一个饱和仿射开子集 \(W'\),
而它包含于仿射开子集 \(V'\cdot x\).
于是 \(W'\) 在 \(X\) 中的像包含 \(x_1,\ldots,x_n\),
由第 V 篇, 4.1(ii), 它是 \(X\) 的仿射开子集.
